% Section fragment for the Next manuscript.
% Requires graphicx; figure files are supplied separately from archived data.
% Single-column figures; run-level statistics are specified in each caption.

\section{Implementation and Evaluation}
\label{sec:evaluation}

Next's evaluation connects recipient-side continuation to funded closure
and identity-preserving repair. We examine dependent payments first,
then source failures, repair costs, organization boundaries, and operating
capacity. These measurements exercise the mechanisms of
Section~\ref{sec:protocol} under the model in Section~\ref{sec:model};
Section~\ref{sec:security} establishes their safety properties under the stated model.

\noindent\textbf{Implementation and measurement.}
The Go implementation uses atomic state overlays, independent block
followers, and CometBFT v0.38.26 modified to enforce complete block
identities and authorized historical installation. Threshold RSA
adaptation uses CIRCL v1.6.5 with trusted offline dealer initialization.
Network experiments run on one Mac Studio with an M4 Max, 16 CPU cores,
64\,GiB RAM, macOS~26.6.2, and Go~1.27.1. Four committee processes and,
per organization, four members and one gateway yield nine or fourteen
service processes. Latency/capacity trials primarily use in-memory
member and committee backends and bbolt wallets with NoSync; signatures
and atomic state updates remain enabled. Source-repair and batch-integration
trials use disk-backed committee application state and BlockStore.
Each run starts from fresh genesis. Current continuation and recovery trials use final user FUEL.
Storage, paired-background, and short-throughput controls use authorized
organization-paid FUEL.
The archived continuation, operating-range, and batch-integration series
use authorized fee sponsorship. The supplement records the build
and settings for each series; batch integration is reported separately
from the earlier structural measurements.

Certified-receipt latency generally starts at the first actual wallet
HTTP send and ends after recipient verification and atomic local recording.
Capital/source-failure drivers include a small amount of pre-send
construction and are compared within their own series. Public and member
completion are observation endpoints, not internal commit timestamps.
Completion throughput includes final background drain. Cross-site ACK
timing uses the sender's monotonic clock and includes the return path.

\subsection{Continuous Spending Without Per-Hop Settlement}
\label{sec:eval-continuation}

A 100-hop payment chain reaches its last certified receipt in a median
158.846\,ms, compared with 65.204\,s when successive hops wait for public
confirmation. Three wallets transfer 100 CAL in rotation, constructing
each payment online from its predecessor's output and paying fees from
their own final FUEL. Three repetitions per mode produce six runs and
600 payments on the nine-service in-memory configuration. Both modes
end at the last certified receipt; the waiting baseline inserts public
confirmation waits \emph{between} hops. Committee commit, flush, and
gossip settings are 250, 10, and 10\,ms, respectively. Under these
settings, public execution and its observation cost about 0.65\,s per
hop in the waiting baseline; certified continuation removes this
inter-hop wait.

Figure~\ref{fig:continuation} reports the six whole-chain durations.
The fast runs take 158.492--159.905\,ms. All 297 fast-mode successor
requests consume the preceding certified output and are sent after its
certified receipt but before its public completion is observed. A
one-input, one-output TXCer is 779 bytes; requests are 1,979 bytes at the
first hop and 2,766 bytes thereafter. For this fixed transaction shape,
direct evidence size stays constant across all 100 hops. These observations
show online dependent spending without inserting a public-confirmation
wait on the recipient's critical path.

\begin{figure}[t]
  \centering
  \includegraphics[width=\columnwidth]{continuation-current.pdf}
  \caption{Time through the last certified receipt of a 100-hop chain.
    Dots are three independent runs per mode; horizontal bars mark their
    medians. The logarithmic time axis includes the same endpoint in
    both modes. Only the waiting baseline requires public confirmation
    between successive payments.}
  \label{fig:continuation}
\end{figure}

\noindent\textbf{Synchronous member commits.}
Three paired 100-hop trials use a nine-service disk-backed configuration,
NoSync wallets, user-paid FUEL, and the same 250/10/10-ms consensus
settings. Each pair differs only in member storage (NoSync or synchronous;
order N/S, S/N, N/S), followed by a synchronous waiting chain.
Median durations are 681.127\,ms for NoSync continuation, 4.471\,s for
synchronous continuation, and 65.013\,s for synchronous waiting.
Thus certified continuation remains 14.5 times faster than waiting
under synchronous member commits. All 900 payments complete with
agreeing committee states; the supplement reports each run.

On the archived delivery-gate build, removing the post-certification delivery gate also shortens continuation
under matched offered load. The gated variant waits for three members to
install complete material or for authenticated public success, whichever
occurs first; it collects no second signature quorum. At 200 payments/s
for 30\,s, three paired runs per mode each send and complete 6,000 payments.
The median of run-level receipt P50s is 1.300\,ms with immediate delivery
and 2.167\,ms with the gate, a 40.0\% reduction. Separate 100-hop trials
reduce median chain time from 287.983 to 200.356\,ms, or 30.43\%
(Fig.~\ref{fig:delivery-gate}). The gain removes a replication/completion
wait while preserving the same three-signature authorization.

\begin{figure*}[!t]
  \centering
  \includegraphics[width=\textwidth]{delivery-gate.pdf}
  \caption{Delivery-gate effect (ms). Left: run-level receipt P50 at
    200 payments/s, 6,000 payments/run, three paired repetitions.
    Right: whole-chain time for three 100-hop chains per mode. Columns
    show the median of the three run values; dots show individual runs.}
  \label{fig:delivery-gate}
\end{figure*}

\subsection{Funded Closure and Source Recovery}
\label{sec:eval-liabilities}

\noindent\textbf{Source-recovery matrix.}
Nine runs use seeds 23, 37, and 59 at withheld-parent rates of 0\%, 1\%,
and 5\%. Each run offers 20 two-hop units/s for 60\,s (1,200 units,
2,400 payments) against a fixed 60,000-CAL reserve with user-paid FUEL
and no top-up.
The nine-service deployment uses disk-backed committee application state
and BlockStore, in-memory members and gateway, and NoSync wallets.
Withheld parents are released after their compensation is committed and
installed on all four replicas, so each withheld source is compensated
after the 30-s public-time deadline and executes later.

All 21,600 payments close. The 216 compensations (12 per 1\% run, 60 per
5\% run) are each followed by recovery. At the end of every run no
obligation is Open or Repaired, the gap is zero, net Spent and Reserved
are zero, and the reserve is back at 60,000\,CAL
(Table~\ref{tab:eval-audit}). Figure~\ref{fig:reserve-recovery} shows the three 5\% trajectories. The lowest sampled
reserve is 59,800--59,900\,CAL at 1\% and 59,600\,CAL at 5\%; it reflects
the driver's release schedule, not a reserve-sizing result. Run-level
receipt P50s are 1.185--1.251\,ms for normal child payments and
1.171--1.207 and 1.176--1.190\,ms for withheld parents and their
children. After the deadline, the committee commit is observed after a
median of 1.003--1.254\,s and installation on all four replicas after
1.257--1.761\,s (P95 up to 2.015\,s), including 250-ms observation
polling. Each normal payment spends 94 FUEL and each repair adds 5;
rewards, burn, and refunds sum to the fee caps with zero Held.

\begin{table}[t]
  \caption{Closing audit of the source-recovery matrix. Each column is
    reproduced by all three seeds; each run has 2,400 payments. Reserve
    values are CAL, fee values are in thousands of FUEL, and fee caps sum the
    per-payment maxima.}
  \label{tab:eval-audit}
  \centering
  \small
  \begin{tabular}{@{}lrrr@{}}
    \toprule
    Withheld parents & 0\% & 1\% & 5\% \\
    \midrule
    Repairs / recoveries & 0 & 12 & 60 \\
    Paid / recovered CAL & 0 & 1,200 & 6,000 \\
    Final reserve & 60,000 & 60,000 & 60,000 \\
    Unclosed / net Spent & 0 & 0 & 0 \\
    Total fee caps & 2,400 & 2,400 & 2,400 \\
    Rewards & 201.600 & 201.660 & 201.900 \\
    Burned & 24.000 & 24.000 & 24.000 \\
    Refunded & 2,174.400 & 2,174.340 & 2,174.100 \\
    Remaining Held & 0 & 0 & 0 \\
    \bottomrule
  \end{tabular}
\end{table}

Ledger-level regression covers the arrival orders directly: all 24
arrival orders of a four-hop chain, with and without due compensation,
under three issuer sequences (one organization, alternating two, and
three with a repeated first issuer) give 144 trajectories. A further 48 split--merge trajectories vary all arrival orders of a
40/60 split, two branches, and their merge under two compensation modes.
Every step in the 192 trajectories conserves the independently computed asset projection, and every
trajectory ends with each issuer's reserve restored, Spent and Reserved
zero, $\mathsf{Paid}=\mathsf{Rec}$, and the consumed inputs still consumed.

\begin{figure}[t]
  \centering
  \includegraphics[width=\columnwidth]{recovery-reserve.pdf}
  \caption{Sampled reserve trajectories for the three 5\% runs (60
    compensated sources each). Recovery restores the initial 60,000 CAL.
    The vertical axis is restricted to 59,600--60,000 CAL to show temporary
    use. Samples reflect the driver's release schedule.}
  \label{fig:reserve-recovery}
\end{figure}

Two additional fourteen-service controls distinguish permanent loss
from delayed recovery. When the parent never arrives, the grandchild
retains its 100 CAL while the responsible issuer's reserve remains
100 CAL lower. When a cross-organization chain arrives grandchild first,
the intermediate source closes its own direct obligation; the root
issuer's compensated 100 CAL is recovered when the root later executes.
The four committee states agree in both cases. An absent parent's own
fee and work reservations can remain outstanding.

\noindent\textbf{Paired end-state check and interference.}
Three pairs of fourteen-service runs, ordered normal/repair,
repair/normal, and normal/repair, use the initial source-recovery
build, before the batch-result tag separation, with committee and member databases
on disk and member and wallet NoSync, each send 4,000 payments at 100
payments/s for 40\,s. Alongside them, four additional two-hop demonstrations have
sources that are delivered normally (normal run) or withheld,
compensated after 30\,s, and later recovered (repair run). All 4,000
payments complete in all six runs. The final CAL held by each owner role
(30,000 and 830,000) and both reserve balances ($10^{12}$) coincide across
the six runs, and the repair runs record 400\,CAL paid and 400\,CAL
recovered. Receipt P50s for normal and repair runs are 1.070/1.127,
1.089/1.108, and 1.083/1.097\,ms, and P95s are 26.84/59.23, 29.49/58.44,
and 27.41/55.15\,ms; completion rates are 97.2--97.8 payments/s in all
runs. The final allocation is consistent with Theorem~\ref{thm:security-neutrality}.
Receipt P50 is 1.3--5.3\% higher and P95 is 1.98--2.21 times the paired
normal values. Concurrent repair therefore affects the receipt tail in
this configuration; these pairs do not isolate the responsible stage.

\noindent\textbf{Storage control.}
Six runs of 1,000 payments at 100 payments/s on the nine-service
configuration (committee on disk, wallets NoSync) differ only in the
member store. With NoSync members, receipt P50 is 1.002--1.035\,ms and
P95 33.97--40.62\,ms, with 93.11--93.93 payments/s completed including
drain. With synchronous member commits, P50 is 43.19--51.85\,ms and P95
80.98--96.56\,ms, with 90.99--92.72 payments/s. All payments complete.
This isolates the cost of synchronous member commits; wallets remain
NoSync, so this is a member-commit control under the execution model of
Section~\ref{sec:model}.

Real external capital sustains longer operation in the earlier-build
capital series, in which no compensation occurred. Separate 300-s
experiments use one Worker per member and ample user FUEL, with one run
per condition. Starting at 3,600 CAL, constant-rate, rate-step, and
3-s delayed-parent workloads each complete 6,000 two-hop units through the
transfers in Table~\ref{tab:eval-capital}. The controller uses the
experiment's declared offered flow. Corresponding fixed-high-reserve
controls complete; a fixed 3,600-CAL control closes 46 of 6,000 units and
retains partial approvals. The supplement gives the controller and all
seven conditions, including their admission limits.

\begin{table}[t]
  \caption{Actual capital additions. Each condition has one 300-s run,
    starts with 3,600 CAL, and completes 12,000 payments. Both columns
    are CAL; every addition debits an external funded account.}
  \label{tab:eval-capital}
  \centering
  \small
  \begin{tabular}{@{}lrr@{}}
    \toprule
    Workload & Final reserve & Added capital \\
    \midrule
    Constant: 20 units/s & 14,400 & 10,800 \\
    Step: 10 to 30 units/s & 21,600 & 18,000 \\
    Parent delayed by 3\,s & 23,300 & 19,700 \\
    \bottomrule
  \end{tabular}
\end{table}

\subsection{Sharing Repair Work While Preserving Successors}
\label{sec:eval-repair}

Same-block batching shares representation work without merging the
identities of the compensated obligations. Structural microbenchmarks on
Windows amd64 with Go~1.27.1 place $k\in\{1,2,4,8\}$ missing inputs in one
historical block and part. Across five repeats, single repairs require
$k$ commands and $k$ part adaptations, whereas a batch uses one of each;
input adaptations remain $k$. At $k=8$, total encoded command bytes fall
from 127,336 to 2,752 (Fig.~\ref{fig:repair-cost}).

The dispersed control separates batching from compact encoding. Eight
independent target blocks still require eight commands and eight part
adaptations, while bytes fall from 30,008 to 5,440. At $k=1$, both formats
perform one part adaptation and encode 3,751 versus 680 bytes. Thus,
compact fields explain part of the byte reduction, while shared block
representation explains the reduced adaptation count. These are encoded
command sizes, not total network traffic or total cryptographic work.

\begin{figure*}[!t]
  \centering
  \includegraphics[width=\textwidth]{repair-cost.pdf}
  \caption{Same-block, same-part repair costs versus obligation count
    $k$. Panels show command/part-adaptation counts and total encoded
    command payloads in KiB (log scale). Five repeats per configuration give
    identical structural measurements. Input adaptations remain $k$.}
  \label{fig:repair-cost}
\end{figure*}

Cross-layer regression executions check successor preservation and
economic consistency. For a payment using 40/60 inputs and its already
executed successor, source repair preserves original authorizations,
output identities, and complete block identities.
Comparisons at a fixed prestate, time, and item order yield the same
economic projection as sequential compensation. Replaying the same
original public history reproduces each height's application hash.
The regressions also exercise cross-part faults, atomic batch rejection,
duplicate commands, and interleaved-target installation.

An archived fourteen-service, disk-backed batch-integration run combines four actual
missing sources into a single 1,568-byte batch whose BatchID agrees
across all committee replicas. After late-parent handling, all 1,000
subsequent payments offered at 100 payments/s complete. Receipt P50/P95
are 1.048/26.531\,ms; total completion time is 11.120\,s, including
observation and shutdown. Gaps and outboxes drain, and committee state
hashes agree. The structural measurements isolate representation costs;
this integration run exercises repair and subsequent payment processing
across the deployed services.

\begin{figure}[t]
  \centering
  \includegraphics[width=\columnwidth]{network.pdf}
  \caption{ACK latency (ms) versus injected cross-site HTTP RTT (ms).
    Each point is the median of three run-level P50s for the first
    120-s cohort; error bars span those P50s' minimum and maximum.
    Aggregate offered rate is 100 payments/s. Committee P2P is undelayed.}
  \label{fig:network}
\end{figure}



\subsection{Continuation Across Registered Organizations}
\label{sec:eval-network}

Cross-organization continuation follows the same certificate path, with
latency reflecting the injected communication distance. On an archived
build recorded in the supplement, the two-organization suite completes 342,000 payments over 24 runs. Two wallet processes each
hold 32 identities across 64 lanes; each organization supplies 50 payments/s.
The latency comparison uses ten-hop same/cross-organization workloads at
0, 20, and 100\,ms injected HTTP RTT, three repetitions per condition.
Each proxy direction adds half the RTT; intra-organization certification
and committee P2P receive no delay injection.

For the common first-120-s cohort, the median of three run-level cross-group
ACK P50s is 1.632, 21.658, and 101.967\,ms, respectively; same-group values
span 1.615--1.744\,ms (Fig.~\ref{fig:network}). ACKs include the return path,
but a recipient may initiate its successor upon acceptance without waiting
for that ACK to reach the previous sender. At 100-ms cross-group RTT,
three 300-s runs each complete 30,000 payments; pending peaks are
106--110 and drain without sustained tail growth. The experiment isolates HTTP path delay within one physical host.



\subsection{Operating Capacity and Service Availability}
\label{sec:eval-operating}

\noindent\textbf{Current-build independent-input throughput.}
Three 30,000-payment runs target 2,000 payments/s on the nine-service
in-memory configuration with NoSync wallets and authorized organization-paid
FUEL. Member processes use a 200\% Go GC target and 16 runtime processors.
All payments complete and the four committee states agree. Completion
rates including drain are 1,750.27, 1,820.91, and 1,887.31 payments/s
over 15.9--17.1\,s, with receipt P50s of 2.646, 2.842, and 2.824\,ms and
P95s of 47.9, 49.1, and 48.0\,ms. The driver recorded
6,551--7,570 payments that had to wait for a fast-receipt admission slot and the P95 dispatch lag was 179--229\,ms, so the
offered load was not uniform at 2,000 payments/s. These are short-run completion measurements; the supplement separately
identifies the build and duration of the 150,000-payment series.


On archived builds, redundant certification and installed evidence
sustain service under the
specified failures. Nine runs at 200 payments/s complete 216,000 payments
with one member delayed or paused for 30\,s. Gateway-loss controls with
one complete installed copy show backup submission 2.503--2.508\,s after
INSTALL; progress requires a surviving complete copy. Conflict controls
produce neither dual conflicting QCs nor repeated principal/fee effects,
while partial approvals remain a capacity cost.

The supplement retains frozen-build mappings and complete per-run
records, including multi-identity CPU/HTTP costs and historical memory
growth separately from backlog. Its 2,048 identities share clients rather
than representing 2,048 processes. It also documents a separately configured LND deployment reference with its own completion endpoint and persistence settings.
